%! Properties and Applications of Determinants — Unit 5, Lesson 5.2 — Student Packet Cover Sheet
\documentclass[10pt]{article}
\usepackage{linalg-article}
\usepackage{linalg-boxes}
\usepackage{ltablex}
\keepXColumns

\begin{document}

% Full-bleed burgundy banner
\begin{tikzpicture}[remember picture, overlay]
  \fill[burgundy] (current page.north west)
    rectangle ([yshift=-0.9in]current page.north east);
\end{tikzpicture}
\vspace{-0.8in}

\begin{center}
  {\color{white}\textbf{\LARGE Linear Algebra}}\\[4pt]
  {\color{blushmid}\large\textbf{Unit 5 \quad Determinants and Linear Transformations}}\\[6pt]
  {\color{white}\textbf{\large Lesson 5.2 \quad Properties and Applications of Determinants}}
\end{center}

\vspace{0.22in}
\namedateperiod
\vspace{0.18in}

\begin{learningtargetbox}
\small
\begin{enumerate}[leftmargin=1.5em, itemsep=5pt]
  \item \ldots{} use the \textbf{rules} of determinants --- $\det I=1$, a \textbf{swap} negates, scaling a row by $t$ scales $\det$ by $t$, and \textbf{adding a multiple of one row to another changes nothing}.
  \item \ldots{} compute a determinant the \textbf{fast} way: reduce to triangular and take $\det=\pm(\text{product of the pivots})$.
  \item \ldots{} use the \textbf{product rule} $\det(AB)=\det A\,\det B$, so $\det A^{-1}=1/\det A$ and $\det A=0$ means \textbf{no inverse}.
\end{enumerate}
\end{learningtargetbox}

\vspace{0.14in}

\begin{tocbox}
\small
\renewcommand{\arraystretch}{1.5}
\begin{tabularx}{\linewidth}{@{} c l X r @{}}
\rowcolor{blush}
\textbf{\#} & \textbf{Component} & \textbf{Description} & \textbf{Score} \\
\midrule
1 & Warm-Up        & Spiral: a $2\times2$ determinant, one elimination step, a triangular determinant & \blank{1.2cm} \\
2 & Guided Notes   & The founding rules; shear invariance; $\det=$ product of pivots; the product rule & \blank{1.2cm} \\
3 & Group Activity & Rules, Pivots, and Products --- computing by elimination and using $\det(AB)$      & \blank{1.2cm} \\
4 & Exit Ticket    & Quick check on determinant by pivots, the product rule, and a justification        & \blank{1.2cm} \\
5 & Homework       & Determinants by elimination, the row rules, products, inverses, and Cramer's rule   & \blank{1.2cm} \\
\midrule
  & \multicolumn{2}{r}{\textbf{Total}} & \blank{1.2cm} \\
\end{tabularx}
\end{tocbox}

\vspace{0.10in}

\begin{remindbox}
\small Expanding a $3\times3$ by hand is real work --- so instead we use the \textbf{rules}
determinants obey. Each rule is really a fact about \emph{volume}: a \textbf{swap} reflects the box
(sign flips), a \textbf{scale} stretches it, and a \textbf{shear} (adding a multiple of one row to
another) slides it without changing base or height, so the volume --- and the determinant --- stay the
same. Because a shear is just an elimination step, you can reduce any matrix to triangular and read
$\det=\pm(\text{product of the pivots})$. And determinants \emph{multiply}: $\det(AB)=\det A\,\det B$.
\end{remindbox}

\end{document}
